% ==============================================================================
% PREFACE.TEX — LỜI NÓI ĐẦU VÀ HƯỚNG DẪN TIẾP CẬN GIÁO TRÌNH
% ==============================================================================

\chapter*{Lời Nói Đầu}
\addcontentsline{toc}{chapter}{Lời Nói Đầu}

Chào mừng bạn đến với \textbf{\booktitle}.

Trong làn sóng bùng nổ của Trí tuệ Nhân tạo Tạo sinh (Generative AI) và các Mô hình Ngôn ngữ Lớn (LLMs), khoảng cách giữa việc \textit{``dùng thử vài câu lệnh prompt''} và việc \textit{``xây dựng một hệ thống AI vận hành an toàn, tin cậy, chịu tải hàng triệu người dùng trong môi trường doanh nghiệp''} là vô cùng lớn.

Hầu hết các tài liệu hiện nay hoặc quá hàn lâm nặng về toán lý thuyết của các nhà nghiên cứu (AI Researchers), hoặc quá hời hợt chỉ dừng lại ở mức hướng dẫn copy-paste code từ các thư viện giao diện cấp cao. Khi hệ thống gặp lỗi ảo giác (\textit{hallucination}), tràn bộ nhớ GPU (\textit{CUDA Out of Memory}), rò rỉ dữ liệu nhạy cảm (\textit{prompt injection}) hoặc sập dịch vụ do độ trễ tăng vọt, người học hoàn toàn bất lực vì thiếu nền tảng kỹ nghệ hệ thống.

\section*{Mục Tiêu Tối Thượng Của Giáo Trình}

Bộ giáo trình này được viết ra với mục đích định hình lại toàn diện phương pháp đào tạo kỹ sư AI:
\begin{enumerate}
  \item \textbf{Khởi đầu từ số 0}: Không giả định người học đã biết trước các thuật ngữ máy tính chuyên sâu. Từng khái niệm đều được dẫn nhập bằng trực giác đời sống và mô hình tư duy trước khi đi vào công thức hoặc mã nguồn.
  \item \textbf{Cốt lõi kỹ nghệ phần mềm}: Bạn sẽ được rèn luyện từ lập trình Python bất đồng bộ, kiến trúc hệ điều hành Linux, mạng máy tính TCP/IP, thiết kế cơ sở dữ liệu quan hệ PostgreSQL và API hiệu năng cao với FastAPI.
  \item \textbf{Hiểu sâu từng bánh răng kiến trúc}: Mổ xẻ toán học của cơ chế Attention, bộ nhớ đệm KV Cache, cách Qdrant tổ chức đồ thị tìm kiếm HNSW, và cơ chế PagedAttention tăng tốc suy luận của vLLM.
  \item \textbf{Bảo mật và an toàn là yếu tố sống còn}: Mỗi giải pháp kỹ thuật đều phải trả lời các câu hỏi: \textit{``Nó có thể bị tấn công như thế nào? Làm sao phát hiện và cô lập rủi ro?''}.
\end{enumerate}

\section*{Phương Pháp Học Tập Tích Cực Với AI}

Giáo trình khuyến khích bạn làm việc cùng các trợ lý AI (AI Pair Programming) theo chu trình kỹ sư chuẩn mực:
\begin{center}
\textbf{Hiểu bản chất} $\rightarrow$ \textbf{Lập kế hoạch} $\rightarrow$ \textbf{Tham vấn AI} $\rightarrow$ \textbf{Hiện thực hóa} $\rightarrow$ \textbf{Kiểm thử} $\rightarrow$ \textbf{Debug} $\rightarrow$ \textbf{Giải thích được}.
\end{center}

Tuyệt đối không áp dụng thói quen thụ động: \textit{Prompt $\rightarrow$ Sao chép code mù quáng $\rightarrow$ Chạy thử $\rightarrow$ Không hiểu nguyên nhân khi xảy ra lỗi}.

Hãy sẵn sàng bước vào hành trình trở thành một \textbf{Kỹ sư AI thực chiến}.
